% Pattern: proof/attach/remark. A remark environment right after a statement,
% and a proof right after the remark: the proof proves the statement. A remark
% after a statement that already has its proof, and an example, keep their own
% proofs.
% Seen in the gold pilot (set E).
% want lemma-1 statement/lemma 1 status=here
% want remark-1 remark 1
% want proof-1 proof - proves=lemma-1
% want lemma-2 statement/lemma 2 status=here
% want proof-2 proof - proves=lemma-2
% want remark-2 remark 2
% want proof-3 proof - proves=remark-2
% want lemma-3 statement/lemma 3 status=none
% want example-1 example 1
% want proof-4 proof - proves=example-1
\documentclass{article}
\newtheorem{lemma}{Lemma}
\newtheorem{remark}{Remark}
\newtheorem{example}{Example}
\begin{document}
\begin{lemma}
A bound.
\end{lemma}
\begin{remark}
The bound is sharp.
\end{remark}
\begin{proof}
A step.
\end{proof}
\begin{lemma}
A helper.
\end{lemma}
\begin{proof}
Another step.
\end{proof}
\begin{remark}
The helper holds in a wider case.
\end{remark}
\begin{proof}
A third step.
\end{proof}
\begin{lemma}
A last one.
\end{lemma}
\begin{example}
A case.
\end{example}
\begin{proof}
A fourth step.
\end{proof}
\end{document}
